\section{Command Line Options}
\label{sec:command_line} 

Several command line options have been added to allow for semi-automated running of tasks or convenient usage. This also allows for (limited) automated batch processing of data. \\

\includegraphics[width=0.5cm]{../../data/icons/hint.png} Please note that configuration of the application still has to be performed using the GUI, therefore it is required to set up the application correctly before command line options can be used successfully.\\

\includegraphics[width=0.5cm]{../../data/icons/hint.png} Please also note that error or warning messages (or related confirmations) will still halt the automatic running of tasks, to ensure that the user is always aware of occurring issues. \\

To get a list of available command line options use e.g.:
\begin{lstlisting}
./COMPASS_1.1.0-x86_64.AppImage --help
\end{lstlisting}

\includegraphics[width=0.5cm]{../../data/icons/hint.png} Please note that the examples in this section use the AppImage. On Windows, \texttt{compass\_client.exe} accepts the same options, e.g.:
\begin{lstlisting}
compass_client.exe --help
\end{lstlisting}

\begin{lstlisting}
Allowed options:
  --help                                shows this help message
  -r [ --reset ]                        resets the user configuration and data
  --override_cfg_path arg               overrides 'default' config subfolder to
                                        other value, e.g. 'org'
  --expert_mode                         sets the expert mode
  --create_db arg                       creates and opens a new DuckDB database
                                        with the given filename, e.g.
                                        '/data/file1.db'
  --open_db arg                         opens an existing DuckDB database with
                                        the given filename, e.g.
                                        '/data/file1.db'
  --import_data_sources_file arg        imports a data sources JSON file with
                                        the given filename, e.g.
                                        '/data/ds1.json'
  --import_view_points arg              imports a view points JSON file with
                                        the given filename, e.g.
                                        '/data/file1.json'
  --import_asterix_file arg             imports an ASTERIX file with the given
                                        filename, e.g. '/data/file1.ff'
  --import_asterix_files arg            imports multiple ASTERIX files with the
                                        given filenames, e.g.
                                        '/data/file1.ff;/data/file2.ff'
  --import_asterix_pcap_file arg        imports an ASTERIX PCAP file with the
                                        given filename, e.g. '/data/file1.pcap'
  --import_asterix_pcap_files arg       imports multiple ASTERIX PCAP files
                                        with the given filenames, e.g.
                                        '/data/file1.pcap;/data/file2.pcap'
  --import_asterix_file_line arg        line to import the ASTERIX data into,
                                        e.g. 'L2'
  --import_asterix_date arg             date of the imported ASTERIX data, in
                                        YYYY-MM-DD format, e.g. '2020-04-20'
  --import_asterix_file_time_offset arg Time of Day offset for the ASTERIX file
                                        import, in HH:MM:SS.ZZZ format, e.g.
                                        '00:00:10.000'
  --import_asterix_ignore_time_jumps    ignores 24h time jumps during the
                                        ASTERIX file import
  --import_asterix_network              imports ASTERIX from the configured
                                        network UDP streams
  --import_asterix_network_time_offset arg
                                        Time of Day offset added during the
                                        ASTERIX network import, in HH:MM:SS.ZZZ
                                        format, e.g. '00:00:10.000'
  --import_asterix_network_max_lines arg
                                        maximum number of lines per data source
                                        during ASTERIX network import, 1..4
  --import_asterix_network_ignore_future_ts
                                        keeps target reports with future
                                        timestamps during the ASTERIX network
                                        import, skips the future timestamp
                                        check
  --import_asterix_network_replay_file arg
                                        simulates live input by replaying the
                                        given IOSS-framed ASTERIX recordings
                                        simultaneously via UDP to a configured
                                        network line, e.g. '/data/file1.ff;/dat
                                        a/file2.ff'
  --import_asterix_network_replay_speed arg
                                        replay speed factor during ASTERIX
                                        network replay, 1.0 replays in recorded
                                        time
  --import_asterix_network_replay_line arg
                                        network line the replayed data is sent
                                        to, L1..L4
  --import_asterix_network_replay_stop_at_end
                                        stops the ASTERIX network import when
                                        the replay file ends
  --asterix_framing arg                 sets ASTERIX framing, e.g. 'none',
                                        'ioss', 'ioss_seq', 'rff', the
                                        configured value is used if not set
  --asterix_decoder_cfg arg             sets the ASTERIX decoder configuration
                                        as JSON string, e.g.
                                        ''{"10":{"edition":"0.31"}}''
                                        (including one pair of single quotes)
  --import_asterix_parameters arg       ASTERIX import parameters as JSON
                                        string, e.g. ''{"filter_modec_active":
                                        true,"filter_modec_max":
                                        50000.0,"filter_modec_min": -10000.0}''
                                        (including one pair of single quotes)
  --import_json arg                     imports a JSON file with the given
                                        filename, e.g. '/data/file1.json'
  --import_gps_trail arg                imports a GPS trail NMEA file with the
                                        given filename, e.g. '/data/file2.txt'
  --import_gps_parameters arg           GPS trail import parameters as JSON
                                        string, e.g. ''{"callsign": "ENTRPRSE",
                                        "ds_name": "GPS Trail", "ds_sac": 0,
                                        "ds_sic": 0, "mode_3a_code": 961,
                                        "set_callsign": true,
                                        "set_mode_3a_code": true,
                                        "set_target_address": true,
                                        "target_address": 16702992,
                                        "tod_offset": 0.0}'' (including one
                                        pair of single quotes)
  --import_sectors_json arg             imports an exported sectors JSON file
                                        with the given filename, e.g.
                                        '/data/sectors.json'
  --list_contexts                       lists all available data contexts and
                                        exits
  --set_context arg                     sets the active data context, e.g.
                                        'Test'
  --calculate_radar_plot_positions      runs the Calculate Radar Plot Positions
                                        task
  --calculate_artas_tr_usage            associates target reports based on the
                                        ARTAS track usage
  --reconstruct_references              runs the Reconstruct References task
                                        (references from sensor and tracker
                                        data)
  --reconstruct_references_cfg arg      reconstructor configuration as JSON
                                        string, e.g. ''{"current_reconstructor_
                                        str": "Scoring + UMKalman"}''
  --analyze_data_source arg             analyzes the data sources of the given
                                        DSType, 'MLAT', 'ADSB' or 'SMR'
  --analyze_parameters arg              analyze data source parameters as JSON
                                        string, e.g. ''{"line_id_tst": 0,
                                        "cell_size_m": 50.0,
                                        "inspector_settings":
                                        {"MLATCoverageInspector":
                                        {"pd_method_int": 1}}}'' (including one
                                        pair of single quotes)
  --load_data                           loads the data into the Views after the
                                        start
  --export_view_points_report arg       exports the view points report after
                                        the start to the given filename, e.g.
                                        '/data/db2/report.tex'
  --evaluate                            runs the evaluation
  --evaluation_parameters arg           evaluation parameters as JSON string,
                                        e.g. ''{"current_standard": "test",
                                        "dbcontent_name_ref": "CAT062",
                                        "dbcontent_name_tst": "CAT020"}''
                                        (including one pair of single quotes)
  --evaluate_run_filter                 runs the evaluation filter before the
                                        evaluation
  --export_report arg                   name of the report to export, e.g.
                                        'EUROCAE ED-87E Evaluation', as PDF if
                                        no export mode is given
  --export_report_directory arg         export directory, e.g. '/data/report2/'
  --export_report_mode arg              export mode, one of 'DocX', 'JSON',
                                        'Latex', 'PDF'
  --no_cfg_save                         does not save the configuration on quit
  --open_rt_cmd_port                    opens the runtime command port (default
                                        27960)
  --enable_event_log                    collects warnings and errors in the
                                        event log
  --quit                                quits after all previous steps are
                                        finished
  --no_highdpi                          disables Qt high-DPI scaling (fixes
                                        double-sized UI on some systems)
  --bundled_font arg                    uses the bundled Ubuntu font at the
                                        given point size, e.g. '11' for
                                        reproducible screenshots
\end{lstlisting}
\ \\

If additional command line options are wanted please contact the author.

\subsection{Options}

\subsubsection{-{}-reset}

Triggers a reset of the user configuration and data. The bundled factory configuration and data are copied into the per-version subfolder of \texttt{\textasciitilde/.compass/}, replacing any user customizations. The short form \texttt{-r} is also accepted.

\subsubsection{-{}-override\_cfg\_path arg}

Overrides the 'default' configuration subfolder name with the given value. Configuration is then read from and written to \texttt{\textasciitilde/.compass/<version>/<arg>/} instead of \texttt{\textasciitilde/.compass/<version>/default/}, which allows separate configuration sets to coexist (e.g. 'org' for organization-specific settings).

\begin{lstlisting}
./COMPASS_1.1.0-x86_64.AppImage --override_cfg_path org
\end{lstlisting}

\subsubsection{-{}-list\_contexts}

Lists all available data contexts and indicates which one is currently active. The application exits immediately after printing the list.

\begin{lstlisting}
./COMPASS_1.1.0-x86_64.AppImage --list_contexts
\end{lstlisting}

Example output:
\begin{lstlisting}
Available contexts:
  Austria  (active)
  Germany
  Test
\end{lstlisting}

\subsubsection{-{}-set\_context name}

Sets the active data context to the given name before any database or import operations are performed. If the context does not exist, the application exits with an error.

\begin{lstlisting}
./COMPASS_1.1.0-x86_64.AppImage --set_context Test --open_db /data/file.db
\end{lstlisting}

\subsubsection{-{}-create\_db filename}

Adds the supplied filename and creates a new DuckDB database.
 
\subsubsection{-{}-open\_db filename}

Adds the supplied filename and opens an existing DuckDB database.

\subsubsection{-{}-import\_data\_sources\_file filename}

Adds the data sources defined in the JSON file to the configuration, as described in \nameref{sec:config_ds_export}.

\subsubsection{-{}-import\_view\_points filename}

After a database was opened, adds the supplied filename and starts an import using the task described in \nameref{sec:ui_import_viewpoints}.
 
\subsubsection{-{}-import\_asterix\_file filename}
\subsubsection{-{}-import\_asterix\_files filenames}
\subsubsection{-{}-import\_asterix\_pcap\_file filename}
\subsubsection{-{}-import\_asterix\_pcap\_files filenames}

After a database was opened, adds the supplied filename and starts an import using the task described in \nameref{sec:ui_import_asterix}.

\subsubsection{-{}-import\_asterix\_file\_line arg}

If an import using the task described in \nameref{sec:ui_import_asterix} is started, it will use the given line identifier (L1 ...L4).


\subsubsection{-{}-import\_asterix\_date arg}

If an import using the task described in \nameref{sec:ui_import_asterix} is started, it will use the import date in YYYY-mm-DD format.

\subsubsection{-{}-import\_asterix\_file\_time\_offset arg}

If an import using the task described in \nameref{sec:ui_import_asterix} is started, the given Time of Day override offset will be applied to all decoded target reports. The argument is in HH:MM:SS.ZZZ format.

\subsubsection{-{}-import\_asterix\_ignore\_time\_jumps}

If an import using the task described in \nameref{sec:ui_import_asterix} is started, 24 hour Time of Day wrap-arounds detected during decoding will be ignored. By default such jumps trigger a day rollover.

\subsubsection{-{}-import\_asterix\_network}

After a database was opened, an import using the task described in \nameref{sec:ui_import_asterix_network} is started.

\subsubsection{-{}-import\_asterix\_network\_time\_offset arg}

If an import using the task described in \nameref{sec:ui_import_asterix_network} is started, the given Time of Day offset is added, in HH:MM:SS.ZZZ format.

\subsubsection{-{}-import\_asterix\_network\_max\_lines arg}

If an import using the task described in \nameref{sec:ui_import_asterix_network} is started, it will use the given maximum number of input lines (and deactivate the others).

\subsubsection{-{}-import\_asterix\_network\_ignore\_future\_ts}

If an import using the task described in \nameref{sec:ui_import_asterix_network} is started, target reports with timestamps in the future relative to the current system time are kept: the future timestamp check is skipped.

\subsubsection{-{}-asterix\_framing framing}

When an Import ASTERIX Task is started the given framing is used, the following options exist:

\begin{itemize}
\item none:  Raw, netto, unframed ASTERIX data blocks, equivalent to the 'raw/netto' value in the GUI
\item ioss:  IOSS Final Format
\item ioss\_seq: IOSS Final Format with sequence numbers
\item rff: Comsoft RFF format
\end{itemize}

\subsubsection{-{}-asterix\_decoder\_cfg 'str'}

When an Import ASTERIX task is started the given configuration is used, in which the editions and mapping can be specified for each category using a JSON string. \\

Using the following string the edition 0.31 can be set for category 010:  \\
'\{"10":\{"edition":"0.31"\}\}' (including one pair of single quotes) \\

In a nice formatting the string looks like this:
\begin{lstlisting}[basicstyle=\small\ttfamily]
'{
"10":
    {
        "edition":"0.31"
    }
}'
\end{lstlisting}
\ \\

Please note the string "10" to identify category 010. \\

For one or a number of categories, the following options can be set:

\begin{itemize}
\item "edition":  ASTERIX editions as string, e.g. "1.0"
\item "ref\_edition":  ASTERIX reserved expansion field as string, e.g. "1.9"
\item "spf\_edition": ASTERIX special purpose field as string, e.g. "ARTAS"
\item "mapping": Mapping used after decoding, e.g. "CAT010 to Radar"
\end{itemize}
\ \\

Please note that the naming must be exactly as in the GUI, otherwise the application quits with an error message.

\subsubsection{-{}-import\_asterix\_parameters}
\label{sec:cli_import_asterix_parameters}

When an Import ASTERIX task is started the given configuration overrides are applied, provided as a JSON string. This is the same set of parameters used by the 'config' argument of the ASTERIX import runtime commands; the available parameters, their units, and examples are described in \nameref{sec:rtcommand_asterix_import_config}.


\subsubsection{-{}-import\_json filename}

After a database was opened, adds the supplied filename and starts an import using the task described in \nameref{sec:ui_import_json}.

\subsubsection{-{}-import\_gps\_trail filename}

After a database was opened, adds the supplied filename and starts an import using the task described in \nameref{sec:ui_import_gps}.

\subsubsection{-{}-import\_gps\_parameters 'str'}

When an Import GPS Trail task is started the given configuration is used, in which the data source and secondary parameters for the GPS trail are defined using a JSON string.

In a nice formatting the string looks like this:
\begin{lstlisting}[basicstyle=\small\ttfamily]
'{
  "callsign":"ENTRPRSE",
  "ds_name":"GPS Trail",
  "ds_sac":0,
  "ds_sic":0,
  "mode_3a_code":961,
  "set_callsign":true,
  "set_mode_3a_code":true,
  "set_target_address":true,
  "target_address":16702992,
  "tod_offset":0.0
}'
\end{lstlisting}
\ \\

Please \textbf{note} that both 'mode\_3a\_code' and 'target\_address' must be given as decimal values.

% \subsubsection{-{}-import\_sectors\_json}
% 
% After a database was opened, adds imports a previously exported sectors JSON file using the following tasks:
% 
% \begin{itemize}
%  \item \nameref{sec:task_manage_sectors}
% \end{itemize}


\subsubsection{-{}-import\_sectors\_json  filename}

After a database was opened, adds the sectors defined in supplied filename using the task described in \nameref{sec:ui_configure_sectors_manage}.

\subsubsection{-{}-calculate\_radar\_plot\_positions}

After a database with imported content exists, radar plot positions are calculated using the task described in \nameref{sec:ui_proc_radar_plot_pos}.

\subsubsection{-{}-calculate\_artas\_tr\_usage}

After a database with imported content exists, target report association can be calculated based on ARTAS tracks using the task described in \nameref{sec:ui_associate_tr_artas}.

\subsubsection{-{}-reconstruct\_references}

After a database with imported content exists, reconstruct reference trajectories using the task described in \nameref{sec:reconst}.

\subsubsection{-{}-reconstruct\_references\_cfg 'str'}

Overrides the reconstructor configuration with a JSON snippet, e.g. to switch the reconstructor or change a parameter for one batch run without altering the GUI configuration.

\begin{lstlisting}
./COMPASS_1.1.0-x86_64.AppImage --open_db /data/file.db \
  --reconstruct_references_cfg '{"current_reconstructor_str": "Scoring + UMKalman"}' \
  --reconstruct_references --quit
\end{lstlisting}

\subsubsection{-{}-load\_data}

Triggers a load process after opening the database.

\subsubsection{-{}-export\_view\_points\_report}

After starting the application into the management window, an export View Points as PDF process is triggered as described in
\nameref{sec:view_points_export_pdf}.

The given argument filename defines the report filename, with the report directory as the parent directory of the given filename.

\subsubsection{-{}-evaluate}

After opening a database a pre-configured evaluation run is triggered as described in \nameref{sec:eval}.

\subsubsection{-{}-evaluation\_parameters 'str'}

When an Evaluation task is started the given configuration is used, in which all parameters of the evaluation can be defined using a JSON string.

In a nice formatting such a string can look like this:
\begin{lstlisting}[basicstyle=\small\ttfamily]
'{
  "active_sources_ref":{
    "CAT062":{
      "1234":true
    }
  },
  "active_sources_tst":{
    "CAT021":{
      "2345":true
    }
  },
  "current_standard":"Dubious Targets",
  "dbcontent_name_ref":"CAT062",
  "dbcontent_name_tst":"CAT021",
  "use_grp_in_sector":{
    "Dubious Targets":{
      "SectorName1":{
        "Optional":false
      },
      "SectorName2":{
        "Optional":false
      },
      "SectorName3":{
        "Optional":true
      },
      "SectorName4":{
        "Optional":false
      }
    }
  }
}'
\end{lstlisting}
\ \\

Please note that 'dbcontent\_name\_ref' cannot be set to CAT001 or CAT048; these categories do not provide the maximum X/Y standard deviation required for reference accuracy filtering and are rejected as reference data. \\

The configuration of such parameters is rather complex, and it is recommended to contact the author if detailed information is needed. \\

The parameter 'report\_figure\_background' sets the background of the geographic report figures, see \nameref{sec:eval_report_figure_background}. It replaces the stored set, a missing field gets its default:

\begin{itemize}
\item color: background color, e.g. "\#000000"
\item map\_show, map\_file, map\_opacity (percent): map
\item borders\_show: national borders
\item sectors\_show, hidden\_sector\_layers (names), hidden\_sectors (sector IDs): sectors. Without 'hidden\_sector\_layers' the default sector selection of the task applies.
\end{itemize}
\ \\

\begin{lstlisting}[basicstyle=\small\ttfamily]
--evaluation_parameters '{"report_figure_background": {"borders_show": false, "map_opacity": 30}}'
\end{lstlisting}
\ \\

\subsubsection{-{}-evaluate\_run\_filter}

When an Evaluation task is started an automatic filtering of targets is performed, as configured using the 'Filter UTNs' dialog defined in \nameref{sec:ui_target_filtering}.

\subsubsection{-{}-analyze\_data\_source 'dstype'}

Triggers the Analyze Data Source task for the given DSType: 'MLAT', 'ADSB' or 'SMR', see \nameref{chap:analyze}.

\begin{lstlisting}
./COMPASS_1.1.0-x86_64.AppImage --open_db /data/file.db --analyze_data_source MLAT --quit
\end{lstlisting}

\subsubsection{-{}-analyze\_parameters 'str'}

Provides a JSON parameter block for the Analyze task, analogous to \nameref{sec:command_line}'s '--evaluation\_parameters' for the Evaluation task. The exact JSON schema depends on the chosen DSType and is documented alongside \nameref{chap:analyze}. \\

The parameter 'report\_figure\_background' sets the background of the geographic report figures, with the same fields as for the Evaluation, see \nameref{sec:analyze_report_figure_background}:

\begin{lstlisting}[basicstyle=\small\ttfamily]
--analyze_parameters '{"report_figure_background": {"color": "#202020", "hidden_sector_layers": []}}'
\end{lstlisting}


\subsubsection{-{}-export\_report 'str'}

Existing reports can be exported as described in the section \nameref{sec:report_export}.

The given argument filename defines the report name, which will be exported with the report directory as the database directory  and the report name as subfolder, as PDF (unless other values are set).

\subsubsection{-{}-export\_report\_directory 'str'}

Sets the report directory to the given value when exporting a report as described in the section \nameref{sec:report_export}.

\subsubsection{-{}-export\_report\_mode 'str'}

Sets the report export type to the given value (one of 'DocX', 'JSON', 'Latex', 'PDF') when exporting a report as described in the section \nameref{sec:report_export}.

\subsubsection{-{}-no\_cfg\_save}

When quitting the application later, no configuration changes are saved.

\subsubsection{-{}-open\_rt\_cmd\_port}

Opens the runtime command TCP port (default 27960) so that the running instance can be controlled from external tools, as described in \nameref{sec:rtcommand_interface}.

\subsubsection{-{}-enable\_event\_log}

Enables collection of warnings and errors emitted during the session into an in-memory event log. The event log can then be queried via the 'get\_events' runtime command, as described in \nameref{sec:rtcommand_interface}.

\subsubsection{-{}-expert\_mode}

Enables expert mode for the session. This grants edit access to the Show DBContent dialog (see \nameref{sec:configure_show_content}) and a few other gated controls. The same setting can be made persistent by setting \texttt{"expert\_mode": true} in \texttt{\textasciitilde/.compass/<version>/task.json}.

\subsubsection{-{}-no\_highdpi}

Disables Qt high-DPI scaling. On some systems (e.g. AlmaLinux 9.7 with GNOME) the reported DPI is inflated, causing a double-sized UI. This flag forces 1:1 pixel scaling. The same effect can be achieved by setting the environment variables \texttt{QT\_AUTO\_SCREEN\_SCALE\_FACTOR=0} or \texttt{QT\_SCALE\_FACTOR=1}.

\subsubsection{-{}-quit}

After the other tasks were run, automatically quits the application.

